\section{Background}
\label{sec:background}

% {Dummy citation}~\citep{alberts2025ctrlz}

Large Language Models (LLMs) are very capable in generating code when prompted to do so. Their ability to generate code that looks correct and compiles as expected does not necessarily mean it always follows best practices. LLMs are trained on real code snippets, which might not be entirely semantically correct. A bad habit of novice programmers when writing HTML code is to overuse the \verb|<div>| tag to define containers, instead of other, more semantically meaningful tags like \verb|<section>|, \verb|<article>|, or \verb|<nav>|. An overuse of the \verb|<div>| tag makes the overall codebase less maintainable and also negatively impacts the accessibility of a website for users using screen readers.
% An example of this can be found in Appendix~\ref{sec:appendix-code-listings} (Listing~\ref{lst:divitis}). % NOTE: I commented out the code listing because there is a page limit of 3 for some reason

Research shows that only a small number of samples can poison LLMs of any size~\cite[]{souly2025poisoningattacksllmsrequire}, which means that this bad habit is very likely to have poisoned the training data of LLMs. Code which was affected by this bad habit looks correct, does not result in rendering errors, and a developer that does not know about the accessibility issues might not even notice the problem. You can configurate an LLM to take these best practices into account when generating code, but it is not guaranteed that it will always follow its system prompt's instructions, especially when running local or smaller models. There exist several Application Programming Interfaces (APIs) online that can be used to check the quality of code and adherence to best practices. One example is the W3C Markup Validation Service\footnote{\url{https://validator.w3.org/}}.

We propose to use such APIs to create a guardrail for LLM outputs which consists of an iterative feedback loop that will continously find errors in the code and regenerate the code by passing the errors the API found in a `regeneration prompt'.

In addition to this issue, there is also a need to evaluate the cost-effectiveness of the proposed guardrail. Such guardrails may allow smaller and cheaper LLMs to achieve the same level of output quality as larger models, which could lead to cost savings. Instead of using a large and expensive LLM to generate code, you could use a smaller and cheaper LLM and give it multiple chances to get the code right, which is possible due to such models being much cheaper to run.
\newpage